\documentclass[10pt,a4paper]{amsart}
\usepackage[margin=19mm]{geometry}
\usepackage{amsmath,amssymb,mathtools}
\usepackage[scr=rsfs,cal=euler]{mathalfa}
\usepackage{xfrac}
\usepackage[unicode,hidelinks,bookmarks=false]{hyperref}
\newcommand{\ie}{i.e.\ }
\newcommand{\cf}{cf.\ }
\newcommand{\resp}{respectively\ }
\newcommand{\Subsection}{\subsection}
\providecommand{\nobreakdash}{\nobreak\hbox{-}}
\setlength{\emergencystretch}{2em}
\multlinegap0pt
\usepackage{cancel}
\usepackage{amssymb}
\usepackage{xcolor}
\usepackage{algorithm}\usepackage{algpseudocode}
\usepackage{mathtools}
\usepackage{float}
\usepackage[shortlabels]{enumitem}
\usepackage[normalem]{ulem}
\newtheorem{lemma}{Lemma}
\newtheorem{theorem}{Theorem}
\def\Hat{\widehat}
\title{The Existence and Stability of Generalized Multi-Source Weber Problems}
\author{V. S. T. Long, N. M. Nam, T. H. Tai, C. Tammer}
\date{Source: arXiv:2607.27900v1 (CC BY 4.0)}
\begin{document}
\maketitle

\noindent\emph{Source and licence.} The Existence and Stability of Generalized Multi-Source Weber Problems. Version arXiv:2607.27900v1 (CC BY 4.0), distributed by arXiv under the Creative Commons Attribution 4.0 International licence, http://creativecommons.org/licenses/by/4.0/, as recorded in the arXiv licence metadata for this exact version and shown on the abstract page https://arxiv.org/abs/2607.27900v1. Full licence notice: \texttt{licenses/arxiv-2607.27900-CC-BY-4.0.txt}.\par\smallskip

\noindent\textit{Editorial scope.} Counted unit: the article's theorem L.O.Solution\_3 (source0.tex lines 1092--1102). The article's complete original proof of it is delivered with it (source0.tex lines 1104--1230, one \texttt{proof} environment of 127 lines). No mathematics is written, repaired or re-derived: the statement and the proof are the article's own lines, in the article's own notation, and no retained line is substituted or re-flowed.  Retained uncounted as the source-local context the proof needs: lines 315--318; lines 351--354; lines 434--442; lines 994--998; lines 1002--1010; lines 1083--1085.  Nothing was omitted from any retained span, so the package carries no omission record.  No retained line contains a citation, so the wrapper carries no bibliography and no bibliography entry.  No retained line contains a figure environment, an image-inclusion command or drawing code, so no image is delivered and the package declares no asset dependency.  Editorial formatting only, in addition: an AMS \texttt{amsart} preamble that restates the article's own declarations and macros used by the retained lines, namely the environments \texttt{lemma}, \texttt{theorem} on separate counters, together with the standard packages those lines need.  The retained environments print in this document as Lemma 1, Theorem 1 in order of appearance, whereas the article prints the same items with the numbering of its own full document.  The complete native source of the article as distributed in the arXiv e-print for this exact version is bundled unchanged as \texttt{sources/206372/source0.tex}.
	\begin{equation}
		\label{minimaltime}
		T_F(x,\Omega) = \inf \left\{ t \geq 0 \mid (x + tF) \cap \Omega \neq \emptyset \right\}, \quad x \in \mathbb{R}^d.
	\end{equation}

\begin{equation} 
	\label{Generalized_Multi_Source_Weber_Problem}
	\min \left\{ f_k(x) = \sum_{i\in I} \min_{j\in J_k} T_F(x^j,\Omega^i) \ \bigg| \ x=(x^1,\dots,x^k)\in\mathbb{R}^{dk} \right\}.
\end{equation}

\begin{lemma}\label{Preliminaries_2}
	Consider the minimal time function defined in \eqref{minimaltime}. Then the following properties hold:
	\begin{enumerate}
		\item [{\rm (a)}] $T_F(x,\Omega) = 0$ if and only if $x \in \Omega$.
		\item [{\rm (b)}] If $\Omega$ is convex, then $T_F(\cdot,\Omega)$ is convex. 
		\item [{\rm (c)}] $T_F(\cdot,\Omega)$ is $\|F^\circ\|$-Lipschitz continuous on $\mathbb{R}^d$.
		\item [{\rm (d)}] $T_F(x,y) \leq T_F(x,z)+T_F(z,y)$ for any $x,y,z \in \mathbb{R}^d$.
	\end{enumerate}
\end{lemma}

\begin{equation}
	\label{L.O.Solution_1}
	I(\overline{x}^j)=\big\{\,i\in I \,\big|\, \Omega^i \in A_{\overline{x}^j} \big\} \text{ and } 
	J_i(\overline{x})=\{j\in J_k \mid \Omega^i \in A_{\overline{x}^j}\}.
\end{equation}

\begin{lemma}
	\label{L.O.Solution_2}
	Let $\overline{x}=(\overline{x}^1,\dots,\overline{x}^k)\in \mathbb{R}^{dk}$ be given. Suppose that for every $i\in I$ the index set $J_i(\overline{x})$ defined in \eqref{L.O.Solution_1} is a singleton. Then there exists $\varepsilon>0$ such that for any 
	$x=(x^1,\dots,x^k)$ satisfying $\|x^j-\overline{x}^j\|<\varepsilon$ for all $j \in J_k$, \color{black}the following properties hold:
	\begin{enumerate}
		\item [\rm {(a)}] $J_i(x)=J_i(\overline{x})$ for all $i \in I$.
		\item [\rm {(b)}] $I(\overline{x}^j)=I(x^j)$ and $A_{x^j}=A_{\overline{x}^j}$ for all $j \in J_k$.
	\end{enumerate}
\end{lemma}

\begin{equation}\label{single_source_Weber_problem}
	\min \biggl\{\,f^j_{1}(x)=\sum_{i \in I(\overline{x}^j)} T_F(x,\Omega^i)\,\bigg|\, x\in \mathbb{R}^d \biggr\}.
\end{equation}

\begin{theorem} \label{L.O.Solution_3}
	Let $\overline{x}=(\overline{x}^1,\dots,\overline{x}^k)\in \mathbb{R}^{dk}$, and suppose that $J_i(\overline{x})$ is a singleton for every $i \in I$. Consider the following assertions:
	
	\begin{enumerate}
		\item[\rm {(a)}] For each $j \in J_k$, if $A_{\overline{x}^j}$ is nonempty, then $\overline{x}^j$ is a (global) optimal solution of the generalized single-source Weber problem \eqref{single_source_Weber_problem} defined on the data set $A_{\overline{x}^j}$.
		\item[\rm{(b)}] $\overline{x}=(\overline{x}^1,\dots,\overline{x}^k)$ is a local optimal solution to problem \eqref{Generalized_Multi_Source_Weber_Problem}.
	\end{enumerate}
	
	Then (a) implies (b). The converse implication also holds if, in addition, each $\Omega^i$ is convex for all $i \in I$.
	
\end{theorem}

\begin{proof} We first prove that $(a) \Rightarrow (b)$. Let $\overline{x}=(\overline{x}^1,\dots,\overline{x}^k)\in \mathbb{R}^{dk}$ satisfy the hypotheses. For each $i \in I$, let $j(i)\in J_k$ denote the unique element of $J_i(\overline{x})$. By Lemma \ref{L.O.Solution_2}(a), there exists $\varepsilon>0$ such that for any $x=(x^1,\dots,x^k)\in \mathbb{R}^{dk}$ satisfying $\|x^j-\overline{x}^j\|<\varepsilon$ for all $j \in J_k$, we have $$J_i(x)=J_i(\overline{x})=\{j(i)\}\quad\text{for every }i \in I.$$ It follows that
	
	\begin{equation}
		\label{L.O.Solution_3_1}T_F(x^{j(i)},\Omega^i)=\min_{j \in J_k}T_F(x^j,\Omega^i).
	\end{equation}
	By assumption, for every $j \in J_k$ such that $A_{\overline{x}^j} \neq \emptyset$, we have
	
	\begin{equation}
		\label{L.O.Solution_3_2}
		\sum_{i\in I(\overline{x}^j)}T_F(\overline{x}^j,\Omega^i)\leq\sum_{i\in I(\overline{x}^j)}T_F(x^j,\Omega^i).
	\end{equation}
	We partition the index set $I$ into disjoint subsets $I = \bigcup_{j \in J_k} I(\overline{x}^j)$ and adopt the convention that$$\sum_{i \in I(\overline{x}^j)} T_F(x^j, \Omega^i) = 0 \quad \text{whenever } I(\overline{x}^j) = \emptyset.$$
	Then we have
	\[
	\begin{aligned}
		f_k(x) 
		&= \sum_{i \in I} \min_{j \in J_k} T_F(x^j,\Omega^i) \\
		&= \sum_{i \in I} T_F(x^{j(i)},\Omega^i) \quad (\text{by \eqref{L.O.Solution_3_1}})\\ 
		&= \sum_{j \in J_k} \sum_{i \in I(\overline{x}^j)} T_F(x^{j(i)},\Omega^i) \\
		&= \sum_{j \in J_k} \sum_{i \in I(\overline{x}^j)} T_F(x^j,\Omega^i) \\
		&\geq \sum_{j \in J_k} \sum_{i \in I(\overline{x}^j)} T_F(\overline{x}^{j},\Omega^i) \quad (\text{by \eqref{L.O.Solution_3_2}})\\
		&= \sum_{i \in I} T_F(\overline{x}^{j(i)},\Omega^i) \\
		&= \sum_{i \in I} \min_{j\in J_k}T_F(\overline{x}^j,\Omega^i) \quad (\text{by \eqref{L.O.Solution_3_1}})\\
		&= f_k(\overline{x}).
	\end{aligned}
	\]
	Thus, for every $x=(x^1,\dots,x^k) \in \mathbb{R}^{dk}$ satisfying $\|x^j-\overline{x}^j \|<\varepsilon$ for all $j \in J_k$, we have
	\begin{equation}
		\label{L.O.Solution_3_3}
		f_k(\overline{x}) \leq f_k(x).
	\end{equation}
	If $\|x-\overline{x}\|<\varepsilon$, then
	\[
	\|x^j-\overline{x}^j\|\le \|x-\overline{x}\|<\varepsilon
	\quad \text{for all } j \in J_k.
	\]
	Hence, by \eqref{L.O.Solution_3_3}, we obtain
	\[
	f_k(\overline{x})\le f_k(x)
	\quad \text{for all } x\in\mathbb{B}(\overline{x};\varepsilon).
	\]
	Therefore, $\overline{x}\in S_k^{{\rm loc}}$.
	
	
	Next, we prove that the converse implication $(b)\Rightarrow(a)$ also holds under the additional assumption that each $\Omega^i$ is convex for all $i\in I$. To this end, let $\overline{x}=(\overline{x}^1,\dots,\overline{x}^k)\in \mathbb{R}^{dk}$ be a local optimal solution of problem~\eqref{Generalized_Multi_Source_Weber_Problem}. This means that there exists $\varepsilon_1>0$ such that
	\[
	f_k(\overline{x})\leq f_k(x) \text{ for all } x \in \mathbb{B}(\overline{x},\varepsilon_1),
	\]
	which is equivalent to
	\begin{equation}
		\label{L.O.Solution_3_4}
		f_k(\overline{x})\leq f_k(x) \text{ for all } x\in \mathbb{R}^{dk} \text{ such that } \|x-\overline{x}\|<\varepsilon_1.
	\end{equation}
	Since $J_i(\overline{x})$ is singleton for any $i \in I$, from Lemma \ref{L.O.Solution_2}(b) it follows that there exists $\varepsilon_2>0$ such that for every $x = (x^1,\dots,x^k)$ satisfying
    \[
    \|x^j-\overline{x}^j\|<\varepsilon_2\quad \text{for all }j \in J_k,
    \]
    we have
	\begin{equation}
		\label{L.O.Solution_3_5}
		A_{\overline{x}^j}=A_{x^j} \text{ and } I(\overline{x}^j)=I(x^j).
	\end{equation}
	It is easy to verify that
	\[
	\|x-\overline{x}\|\leq \sqrt{k}\,\max_{j\in J_k}\|x^j-\overline{x}^j\|.
	\]
	Moreover, for any $K>0$, one has
	\[
	\|x^j-\overline{x}^j\|<K \quad \text{for all } j \in J_k
	\Longleftrightarrow
	\max_{j\in J_k}\|x^j-\overline{x}^j\|<K.
	\]
	Set
	\[
	\varepsilon=\min\left\{\frac{\varepsilon_1}{\sqrt{k}},\varepsilon_2\right\}.
	\]
	Then, for any $x=(x^1,\dots,x^k)\in \mathbb{R}^{dk}$ satisfying
	\[
	\|x^j-\overline{x}^j\|<\varepsilon \quad \text{for all } j \in J_k,
	\]
	the vector $x$ satisfies both \eqref{L.O.Solution_3_4} and \eqref{L.O.Solution_3_5}. Thus
	\begin{equation}
		\label{L.O.Solution_3_6}
		A_{\overline{x}^j}=A_{x^j},\quad I(\overline{x}^j)=I(x^j),\quad\text{and}\quad     f_k(\overline{x})\leq f_k(x).
	\end{equation}
	Assume that $A_{\overline{x}^{j_0}} \neq \emptyset$ for some $j_0 \in J_k$. Let $y \in \mathbb{R}^d$ be arbitrary such that
	\[
	\|y-\overline{x}^{j_0}\|<\varepsilon.
	\]
	Define $\Hat{x}=(\Hat{x}^1,\dots,\Hat{x}^k)\in \mathbb{R}^{dk}$ by
    \begin{equation}
    \label{Hat_x_def}
    	\Hat{x}^j=
	\begin{cases}
		y & \text{if } j=j_0,\\
		\overline{x}^j & \text{if } j \in J_k\setminus\{j_0\}.
	\end{cases}
    \end{equation}
    Since $\|\Hat{x}^j-\overline{x}^j\|<\varepsilon$ for all $j \in J_k$, we obtain
	\begin{equation}
		\label{L.O.Solution_3_7}
		\begin{aligned}
			f_k(\overline{x})&=\sum_{i \in I(\overline{x}^{j_0})}\left(\min_{j \in J_k}T_F(\overline{x}^j,\Omega^i) \right)+\sum_{i \not\in I(\overline{x}^{j_0})}\left(\min_{j \in J_k}T_F(\overline{x}^j,\Omega^i) \right)\\
			&=\sum_{i \in I(\overline{x}^{j_0})}T_F(\overline{x}^{j_0},\Omega^i) +\sum_{i \not\in I(\overline{x}^{j_0})}\left(\min_{j \in J_k\setminus \{j_0\}}T_F(\overline{x}^j,\Omega^i) \right),
		\end{aligned}
	\end{equation}
	and
	\begin{equation}
		\label{L.O.Solution_3_8}
		\begin{aligned}
			f_k(\Hat{x})    &=\sum_{i \in I(\Hat{x}^{j_0})}\left(\min_{j \in J_k}T_F(\Hat{x}^j,\Omega^i) \right)+\sum_{i \not\in I(\Hat{x}^{j_0})}\left(\min_{j \in J_k}T_F(\Hat{x}^j,\Omega^i) \right)\\
            &=\sum_{i \in I(\Hat{x}^{j_0})}T_F(\Hat{x}^{j_0},\Omega^i) +\sum_{i \not\in I(\Hat{x}^{j_0})}\left(\min_{j \in J_k\setminus\{j_0\}}T_F(\Hat{x}^j,\Omega^i) \right)\\
			&=\sum_{i \in I(\Hat{x}^{j_0})}T_F(y,\Omega^i)+\sum_{i \not\in I(\Hat{x}^{j_0})}\left(\min_{j \in J_k\setminus\{j_0\}}T_F(\overline{x}^j,\Omega^i) \right) \quad (\text{by \eqref{Hat_x_def}})\\
			&=\sum_{i \in I(\overline{x}^{j_0})}T_F(y,{\Omega^i})+\sum_{i \not\in I(\overline{x}^{j_0})}\left(\min_{j \in J_k \setminus \{j_0\}}T_F(\overline{x}^j,{\Omega^i}) \right) \quad (\text{by \eqref{L.O.Solution_3_6}}).
		\end{aligned}
	\end{equation}
    It also follows from \eqref{L.O.Solution_3_6} that
	\[
	f_k(\overline{x}) \le f_k(\Hat{x}).
	\]
	Combining this with \eqref{L.O.Solution_3_7} and \eqref{L.O.Solution_3_8}, we get
	\[
	\sum_{i \in I(\overline{x}^{j_0})} T_F(\overline{x}^{j_0},\Omega^i) \le \sum_{i \in I(\overline{x}^{j_0})} T_F(y,\Omega^i) 
	\quad \text{for all } y \in \mathbb{R}^d \text{ such that } \|y-\overline{x}^{j_0}\|<\varepsilon.
	\]
	Hence, $\overline{x}^{j_0}$ is a local optimal solution of the single-source Weber problem associated with the data set $A_{\overline{x}^{j_0}}$. Under the additional assumption that each $\Omega^i$ is convex, it follows from Lemma~\ref{Preliminaries_2}(b) that the function $T_F(\cdot,\Omega^i)$ is convex for every $i\in I$. Therefore, the objective function of problem~\eqref{single_source_Weber_problem} is convex, and hence every local optimal solution is also a global one. Consequently, $\overline{x}^{j_0}$ is a global optimal solution of problem~\eqref{single_source_Weber_problem}. This completes the proof.
\end{proof}

\end{document}
